\chapter*{Objective}
\addcontentsline{toc}{chapter}{Objective}

The objective of this handbook is to establish a coherent systems-level foundation for studying embedded artificial intelligence and autonomous robotics. The central premise is that an embedded-AI system is not merely a neural network running on a small processor. It is a closed-loop cyber-physical system that must convert physical phenomena into measurements, representations, estimates, decisions, and actuator commands while satisfying constraints on latency, energy, memory, communication bandwidth, reliability, and safety.

The handbook develops the conceptual vocabulary needed to connect topics that often appear as separate islands: camera and event-camera sensing, LiDAR and radar measurements, modality selection, calibration and synchronization, sensor fusion, feature learning, attention, graph neural networks, spiking neural networks, odometry, simultaneous localization and mapping, FPGA acceleration, compute-in-memory, runtime reconfiguration, and near-sensor intelligence. Each topic is placed within a common \emph{sensing-to-decision} chain so that a new paper can be located quickly by its inputs, outputs, fusion level, embedded constraint, and system role.

A second objective is to support the Week 7 paper review in EE 6900. The six assigned papers are treated as complementary case studies in multimodal perception, adaptive resource management, event-driven processing, state estimation, and hardware--algorithm co-design. The handbook reconstructs the research question of each paper, identifies the gap being addressed, and explains how the papers relate to one another across the complete autonomous-system pipeline.

The intended learning outcome is practical engineering judgment. After studying this document, the reader should be able to distinguish a sensor from a modality, a measurement from a state, feature-level fusion from state-estimation fusion, loose coupling from tight coupling, event-driven computation from spiking computation, model accuracy from end-to-end system performance, and algorithmic efficiency from deployable hardware efficiency. The final goal is a reusable intellectual framework for reading future embedded-AI literature with greater depth, precision, and critical discipline.
